\section{Discussion}
\label{sec:discussion}

\paragraph{Format is a deployment decision.}
Serialisation strategy is not a cosmetic choice. It is a performance variable with
statistically significant, clinically meaningful effect sizes. For Mistral-7B,
choosing Clinical Narrative (Strategy C) over Raw JSON (Strategy A) raises mean F1
from 0.72 to 0.91. That is a 19 percentage point gain on the same model, same hardware,
same data, with no additional training. For Llama-3.1-8B, the gain is smaller but
still significant: 0.92 to 0.95. These gains cost nothing beyond a preprocessing step.
No retraining, no fine-tuning, no prompt engineering beyond the format change itself.
Any team deploying a 7B--8B model for medication reconciliation should default to
Clinical Narrative format unless they have evidence that their specific model and data
distribution behave differently.

At 70B, the story changes. Llama-3.3-70B reaches near-perfect performance on Raw JSON
(F1 = 0.9956) and shows no significant degradation on Clinical Narrative (F1 = 0.9850),
but drops to F1 = 0.8742 on the Chronological Timeline, with 25 patients receiving zero
F1. A model of this scale appears to handle JSON structure natively without needing it
reorganised, but it still struggles with temporal inference across interleaved statuses,
even when it handles all other formats correctly. Format sensitivity does not disappear
at large scale; it shifts. Large models are more robust to structural format changes
but remain vulnerable to tasks that require specific reasoning patterns.

\paragraph{Omission defines the clinical risk profile.}
Omission is the dominant failure mode across all conditions. In all 20 model--strategy
combinations, mean precision exceeds mean recall, meaning models are substantially more
likely to miss an active medication than to fabricate one. Hallucination does still
occur, particularly for smaller models on less structured formats: Phi-3.5-mini on Raw
JSON has mean precision 0.73 (roughly 27\% of its output items are not in the ground
truth) and Mistral-7B on Raw JSON reaches 0.89. These are measurable, non-trivial error
rates. What the data shows is not the absence of hallucination but its
subordination to omission.

Two caveats sharpen this finding. First, under exact string matching, some apparent
hallucinations reflect name formatting variations rather than true fabrication: the
model outputs a real active medication with a suffix such as \textit{(RxNorm: 314076)}
or a truncated bracketed brand name, which counts as both a false positive and a false
negative. Inspection of Llama-3.3-70B's Strategy D errors shows this pattern accounts
for the majority of its 73 precision errors. Our reported precision and recall values
are therefore conservative lower bounds; fuzzy matching would raise both. Second,
the asymmetry itself is consistent across scales and strategies even at 70B, which
suggests it is a property of the task under exact-match evaluation rather than a
small-model artefact.

A model with high precision and imperfect recall produces output that can be treated
as a confirmed active subset, not a complete active list. Downstream reconciliation
workflows should be designed accordingly. Model output flags confident positives;
human review focuses on completeness rather than accuracy checking. That is a
different cognitive task, and potentially a more tractable one. Instead of verifying
every model output item is correct, the clinician's job becomes catching what the
model missed.

\paragraph{Polypharmacy patients are systematically underserved.}
Small models hit a hard capacity ceiling around 7--10 concurrent medications. For
Mistral-7B, this ceiling is stark: recall at 11 active medications falls to 0.24.
For context, roughly 30\% of adults aged 50--80 take five or more prescription
medications \citep{lantz2020polypharmacy}. Patients with 10+ active medications are
rare in the general population, but they are the patients most likely to present at
hospital handoffs and most at risk from reconciliation errors. The small-model capacity
ceiling affects precisely the population where medication reconciliation matters most.

Llama-3.3-70B does not show this ceiling. It maintains near-perfect recall through
16 active medications. The practical implication: for clinical deployments involving
complex, multi-medication patients, model scale may matter more than serialisation
choice. No format of Mistral-7B reaches Llama-3.3-70B's performance on the hardest
patients.

\paragraph{Domain pretraining without instruction tuning is insufficient.}
BioMistral produced F1 = 0.0 on all 4,000 runs assigned to it. Its sibling model,
Mistral-7B (same architecture, same parameter count), produced F1 = 0.9149 on
Strategy C. The domain knowledge was there. The instruction-following capability was
not. BioMistral was further-pretrained on 3B tokens of PubMed Central text starting
from Mistral-7B Instruct v0.1, which should have preserved instruction-following
capability in principle. It did not. The catastrophic forgetting of structured
response generation during continued pretraining rendered it entirely unusable,
regardless of input format.

This aligns with \citet{kim2025medicalhallucination}, who show that general-purpose
models outperform medical-specialised models on structured clinical tasks (76.6\% vs.\
51.3\% hallucination-free) and interpret the gap as a reasoning failure rather than a
knowledge deficit. Our result is more extreme. BioMistral does not underperform; it
produces zero useful output. For clinical NLP practitioners, the lesson is clear:
instruction-following capability is the prerequisite. A general-purpose instruction-tuned
model at the same parameter count will reliably outperform a domain-pretrained base
model on these tasks, regardless of medical vocabulary coverage.
